\chapter{Автоморфизм орбиты и KO6 half-trace}

Предыдущая глава оставила chiral block
\begin{equation}
 B(x,z)=
 \begin{pmatrix}
 x&\overline x\\
 \overline z&z
 \end{pmatrix}.
\end{equation}
Теперь требуется отделить физические параметры от замен базиса,
автоморфизма обмена листов и обязательного KO6-дублирования.

\section{Reality-совместимые замены базиса}

На двух \(J\)-фиксированных диагональных секторах разрешены вещественные
знаки \(\epsilon_X,\epsilon_C=\pm1\). На паре
\(H_{XC}\leftrightarrow H_{CX}\) разрешена фаза
\(e^{i\theta}\) и её сопряжение. Тогда
\begin{equation}
 x\longmapsto
 \epsilon_Xe^{-i\theta}x,
 \qquad
 z\longmapsto
 \epsilon_Ce^{i\theta}z.
\end{equation}
Следовательно,
\begin{equation}
 |x|,\qquad |z|,\qquad xz\ \text{с точностью до знака}
\end{equation}
являются инвариантами орбиты. Непрерывная фаза
\(\arg(xz)\bmod\pi\) не устраняется.

Автоморфизм \(A_F=\mathbb C\oplus\mathbb C\), меняющий листы местами,
действует как \(x\leftrightarrow z\). Он превращает пару модулей в
неупорядоченную, но не фиксирует их значения.

\begin{proposition}[Automorphism-orbit no-go]
Reality-совместимые rephasing и обмен листов не сводят \(D_F\) к
единственной дискретной орбите. После факторизации остаются неупорядоченные
модули \(\{|x|,|z|\}\) и непрерывная фаза
\(\arg(xz)\bmod\pi\).
\end{proposition}

\section{Спектральные инварианты}

Чётное спектральное действие зависит от сингулярных значений \(B\).
Для них
\begin{align}
 \Tr(BB^\dagger)
 &=2\left(|x|^2+|z|^2\right),\\
 \det(BB^\dagger)
 &=4\,\operatorname{Im}(xz)^2.
\end{align}
Поэтому даже спектр сохраняет две непрерывные комбинации: общий
квадратичный масштаб и величину ориентационно-чувствительной фазы.
Алгебра и first-order condition задают форму потенциала, но не выбирают
его минимум без конкретной spectral function и масштаба.

Если наложить голый fixed-point обмен \(x=z\), остаются модуль и фаза
одного комплексного числа. Если же обмен сопровождается условием
\(z=\overline x\), то \(\operatorname{Im}(xz)=0\) и связующий блок теряет
полный ранг. Таким образом, симметрия обмена сама по себе не создаёт
единственную ненулевую CP-орбиту.

\section{Явное KO6-дублирование}

Пусть четырёхсекторный модуль имеет \((H_4,\gamma_4,J_4)\) с
\(J_4\gamma_4=\gamma_4J_4\). Введём
\begin{equation}
 H_8=H_4\oplus H_4^{c},
 \qquad
 \gamma_8=
 \begin{pmatrix}
 \gamma_4&0\\0&-\gamma_4
 \end{pmatrix},
\end{equation}
а \(J_8\) пусть антилинейно обменивает две копии. Тогда
\begin{equation}
 J_8\gamma_8=-\gamma_8J_8.
\end{equation}
Для \(J\)-инвариантного наблюдаемого \(Q\)
\begin{equation}
 \Tr_{H_8}Q=2\Tr_{H_4}Q.
\end{equation}
Обычный полный trace поэтому удваивает ранее выведенные факторные веса и
rank-вклады.

\section{Reality-orbit half-trace}

Единственная нормировка, инвариантная относительно формального добавления
сопряжённой \(J\)-копии, имеет вид
\begin{equation}
 \Tr_{\mathrm{orb}}Q
 =\frac12\Tr_{H_8}Q.
\end{equation}
Она точно восстанавливает недублированные значения:
\begin{equation}
 \Tr_{\mathrm{orb}}p_X=1,
 \qquad
 \Tr_{\mathrm{orb}}p_C=1,
 \qquad
 \Tr_{\mathrm{orb}}P_H=23.
\end{equation}

\begin{proposition}[Условный half-trace pass]
Если parent action определяется следом на \(J\)-орбитах, а не полным
комплексным следом на удвоенном модуле, KO6-дублирование не меняет
факторные нормы и нейтринный rank. Коэффициент \(1/2\) фиксирован размером
reality-орбиты и не является секторным весом.
\end{proposition}

Однако конечная спектральная тройка сама по себе не требует заменять
бозонный spectral trace на orbit trace. В фермионном интеграле устранение
дублирования связано с выбором физического подпространства или
Pfaffian/Majorana prescription, тогда как бозонное спектральное действие
обычно считает полный спектр.

\begin{nogocriterion}[Half-trace origin gap]
Reality-orbit half-trace считается выводом только если он возникает из
одной заранее заданной функциональной меры и применяется ко всем
градуированным Hessian-блокам. Назначение множителя \(1/2\) только для
восстановления прежних нормировок является скрытым глобальным весом.
\end{nogocriterion}

\section{Вердикт}

Автоморфизмы не фиксируют \(x,z\), а KO6 half-trace проходит лишь
условно. Следующий конструктивный гейт должен задать конкретный
спектральный или полиномиальный parent functional для \(D_F\), проверить
его стационарные орбиты и одновременно вывести orbit-trace prescription.
До этого vertex существует, но его масштаб и CP-фаза не предсказаны.

Машинный аудит сохранён в
\texttt{s2t\_v3\_automorphism\_halftrace\_results.json}.